\section{Tokamak: contener el Sol con campos magnéticos}

La sección anterior explicó el principio de la reacción; esta pasa a las rutas tecnológicas principales. La Tierra no puede contener la reacción de fusión con una gravedad enorme como lo hace el Sol, por lo que la humanidad ha explorado sobre todo dos tipos de confinamiento: el confinamiento inercial y el confinamiento magnético. El confinamiento inercial comprime cápsulas de combustible con láseres intensos de duración extremadamente corta, entre otros métodos; es adecuado para estudiar la física de alta densidad de energía y problemas relacionados con armamento, pero no para la generación eléctrica civil estable. El confinamiento magnético, en cambio, usa campos magnéticos para mantener el plasma a alta temperatura dentro de una región del espacio y evitar que toque directamente las paredes sólidas.

El tokamak es la ruta más extendida dentro del confinamiento magnético y la que ha alcanzado los parámetros experimentales más altos. Su idea intuitiva es un campo magnético en forma de dona: el plasma se mueve a lo largo de una trayectoria anular y queda confinado en la configuración que forman en conjunto el campo magnético toroidal y la corriente toroidal. En la entrevista, Yang Zhao (杨钊) subraya que el tokamak se volvió el consenso no porque su nombre suene bien, sino porque durante las últimas décadas ha superado con claridad a las demás rutas de confinamiento magnético en parámetros experimentales.

\begin{figure}[H]
\centering
\includegraphics[width=0.92\textwidth]{tokamak-stack.png}
\caption{Pila de sistemas del tokamak: imanes, cámara de vacío, control del plasma y conversión de energía, acoplados capa por capa. Diagrama conceptual propio, elaborado a partir de la conversación de 00:06:32--00:16:00 y 01:00:50--01:04:36.}
\end{figure}

\begin{knowledgebox}{Lectura de la figura: el tokamak es ingeniería de sistemas, no un solo imán}
La figura recorre desde los imanes superconductores hasta la cámara de vacío, el sistema de calentamiento, el sistema de control y el sistema termohidráulico, y muestra que el tokamak es una máquina en la que se acoplan múltiples sistemas. Los imanes proporcionan el confinamiento, la cámara de vacío proporciona el entorno de operación, el calentamiento y el control llevan el plasma a los parámetros objetivo, y el sistema termohidráulico determina si puede convertirse en una central eléctrica.
\end{knowledgebox}

\subsection{Cobre, superconductores de baja temperatura y superconductores de alta temperatura}

Esta subsección examina cómo el material de los imanes cambia todo el dispositivo. Los primeros tokamaks solían usar imanes de cobre, de ingeniería sencilla; pero el cobre tiene resistencia eléctrica, se calienta al conducir corrientes intensas y solo permite operar en pulsos cortos. Los superconductores de baja temperatura permiten operar durante periodos largos, pero requieren temperaturas extremadamente bajas y complejos sistemas de vacío y de blindaje criogénico; más importante aún, su campo magnético crítico es limitado, de modo que, para alcanzar una ganancia de energía suficiente, el dispositivo tiene que ser muy grande. El ejemplo típico es un megaproyecto internacional como ITER.

El valor central de los superconductores de alta temperatura no es que «la temperatura realmente sea alta», sino que, en condiciones de operación a baja temperatura, soportan campos magnéticos más intensos. Cuanto más intenso es el campo, menor es el tamaño del dispositivo necesario para un mismo desempeño; y cuanto menor es el tamaño, menores son el costo y el plazo de construcción. El juicio de Yang Zhao sobre la ruta es que los superconductores de alta temperatura podrían convertir los grandes dispositivos de la ruta de baja temperatura, del orden de decenas de miles de millones de euros y con plazos de varias décadas, en objetos de ingeniería más adecuados para que una empresa emergente los desarrolle por iteraciones.

\begin{knowledgebox}{Términos clave: tres generaciones de imanes de tokamak}
\begin{tabularx}{\textwidth}{p{0.18\textwidth}p{0.34\textwidth}X}
\toprule
Etapa & Ventajas & Problemas principales \\
\midrule
Imanes de cobre & Baja complejidad de ingeniería; adecuados para los primeros experimentos & Generan mucho calor; solo operan en pulsos cortos. \\
Superconductores de baja temperatura & Operación prolongada; experiencia acumulada con dispositivos como EAST, KSTAR y JT60SA & Sistema criogénico complejo; el límite del campo magnético crítico obliga a dispositivos enormes. \\
Superconductores de alta temperatura & Campo magnético crítico más alto; prometen reducir de forma notable el tamaño del dispositivo y su costo & El procesamiento del material, los esfuerzos estructurales, la integración criogénica y la validación del sistema todavía son muy recientes. \\
\bottomrule
\end{tabularx}
\end{knowledgebox}

\begin{importantbox}{La lógica comercial de los superconductores de alta temperatura}
Los superconductores de alta temperatura no simplifican la reacción de fusión en sí, sino que cambian la economía de ingeniería del dispositivo: un campo magnético más intenso permite un dispositivo más pequeño, y un dispositivo más pequeño trae menor costo, ciclos de iteración más cortos y una ruta de investigación y desarrollo más adecuada para que la impulse una empresa emergente.
\end{importantbox}

\subsection{Resumen de la sección}

El tokamak es la ruta de confinamiento magnético con mayor acumulación experimental hasta ahora. Lo decisivo en el tokamak de superconductores de alta temperatura no es solo el desempeño científico, sino si puede usar un campo magnético más intenso para reducir la escala del dispositivo y el costo por kilowatt-hora.
